
\section{相关理论与技术}

在此输入本章的引言内容，概述本章将要讨论的
理论基础或技术背景。

\subsection{基础概念}

在此输入基础概念的说明。下面是一个代码示例：

\begin{lstlisting}[language=C, caption={代码示例标题}, label={lst:example}]
#include <stdio.h>

int main() {
    // 代码注释说明
    printf("Hello, World!\n");
    return 0;
}
\end{lstlisting}

\subsection{数据类型}

在此输入相关内容的说明。表 \ref{tab:example_data} 
展示了相关数据。

\begin{table}[H]
    \centering
    \caption{表格标题}
    \label{tab:example_data}
    \begin{tabularx}{0.9\textwidth}{l c X}
        \toprule
        类型 & 数值 & 描述 \\
        \midrule
        类型A & 10 & 类型A的详细描述说明 \\
        类型B & 20 & 类型B的详细描述说明 \\
        类型C & 30 & 类型C的详细描述说明 \\
        类型D & 40 & 类型D的详细描述说明 \\
        \bottomrule
    \end{tabularx}
\end{table}

\subsection{本章小结}

在此总结本章讨论的主要内容和关键概念。
